\subsection{Extraneous Solutions} 
When we clear denominators in a rational equation, we multiply both sides by a \makeblank[1in] expression, which we've avoided so far. Why?
\begin{Example}
Solve $\dfrac{3x+2}{x^2-4x-12}=\dfrac{2x+1}{x^2-2x-8}$ by clearing denominators.
\begin{lpic}{}{19}
\foreach \y/\l/\inst in {1/a/Find the LCD,4/b/Multiply both sides by,9/c/Solve the new equation,16/d/Check}
{
	\regtextright{1}{-\y}{\l. };
	\regtextleft{1}{-\y}{\inst};
}
\foreach \x/\exp in {1/x^2-4x-12,\value{cols}/x^2-2x-8}
	\regtextleft{\x}{-2}{$\exp={}$};
\regtextleft{1}{-3}{$\mathrm{LCD}={}$};
\end{lpic}
Wait---what happened?\lbr
If $x=\makeblanksmall$, $\makeblank=\makeblanksmall$, so we accidentally multiplied both sides by \makeblanksmall. $x=\makeblanksmall$ is an \term{extraneous solution}.
\end{Example}
\begin{definition}[Extraneous Solution]
An \term{extraneous solution} to an equation is one which appears through the solving process, but was not an actual solution to the equation. One reason this may occur is accidentally multiplying both sides of the equation by \makeblanksmall.
\end{definition}